\documentclass{beamer}

\usepackage{mstyle-slides}

\usepackage[brazil]{babel}

\title{Arquivo de Exemplo para o Estilo mstyle-slides.sty}
\author{guilhermeivo}
\date{\today}

\begin{document}

\frame[plain]{\titlepage}

\begin{frame}
\frametitle{Fórmula de Machin}

\begin{itemize}
\item A fórmula de Machin foi descoberta por John Machin (1680--1751).
\item Foi usada para calcular $\pi$ com 100 casas decimais.
\item Posteriormente, William Shanks (1812--1882) utilizou essa fórmula para calcular $\pi$ com 707 casas decimais.
\end{itemize}

\medskip

\begin{equation}
\frac{\pi}{4}=4\arctan\frac{1}{5}-\arctan\frac{1}{239}
\end{equation}

\end{frame}

\begin{frame}
\frametitle{Série da Arcotangente}

\begin{itemize}
\item A função $\arctan(x)$ possui uma expansão em \textbf{série de potências}.
\item Essa série converge para $|x| < 1$.
\end{itemize}

\begin{equation}
\arctan x=\sum\limits_{n=0}^{\infty}(-1)^{n}\frac{x^{2n+1}}{2n+1},\quad\lvert x\rvert<1
\end{equation}

\end{frame}

\begin{frame}[allowframebreaks,plain]
\frametitle{Referências}

\nocite{pi-source-book}
\nocite{special-issue}
\nocite{aocc-user-guide}
\nocite{ieee}

\bibliographystyle{plain}
\bibliography{references}

\end{frame}

\end{document}